\documentclass[a4paper,10pt,pdftex,pagesize]{scrartcl}

\usepackage{ucs}
\usepackage[utf8]{inputenc}
\usepackage{mathpazo}
\usepackage[german]{babel}
\usepackage[OT1,T1]{fontenc}
\usepackage[pdftex]{graphicx}
\usepackage{amsmath}
\usepackage[siunitx,european]{circuitikz}
\usepackage[pdftex]{hyperref}

\author{Thomas Maul}
\title{Übung zu Überlagerung von elektrischen Feldern mit Lösungsvorschlag}
\date{22.09.26}

\begin{document}
% \maketitle

\section{Aufgabe 1}
Gegeben sind zwei Ladungen $\text{Q}_1$
und
$\text{Q}_2$ . Mit dem Koordinaten:
\\ $\text{Q}_1 \left( \begin{array}{c} 1\\ 1\end{array} \right), \text{Q}_2 \left( \begin{array}{c} 5\\2                                                  \end{array}
\right)$
\\
Die Ladungen haben die Werte:
$\text{Q}_1 = 1,2\ nC, \text{Q}_2 = -15\ nC$.\\
Am Punkt Pn (siehe \ref{subsecMesspunkte}) befindet sich eine Ladung $\text{Q}_3$ von $33\ nC$.
\subsection{Aufgaben}
\begin{enumerate}
 \item Zeichne je ein Koordinatensystem mit den Punkten Q1, Q2 und Pn. Pn ist je ein Punkt aus der Liste P1 bis P5.
 \item bestimme jeweils den Abstand der Ladungen Q1 und Q2 zur Ladung Pn.
 \item berechne die Feldstärke, die Q1 in Pn verursacht.
 \item berechne die Feldstärke, die Q2 in Pn verursacht.
 \item überlagere die Felder rechnerisch und bestimmen damit das resultierende Feld.
 \item bestimme die Kraft von Q1 auf Pn.
 \item bestimme die Kraft von Q2 auf Pn.
 \item bestimme die resultierende Kraft von Q1 und Q2 auf Pn mit Betrag und Winkel gegenüber der X-Achse.
 \item bestimme die resultierende Kraft von Q1 und Q2 auf Pn in x-Richtung und in y-Richtung.
\end{enumerate}

Alle Längenangaben sind in Meter.

\subsection{Messpunkte}
\label{subsecMesspunkte}
Für Pn soll jeweils einer folgenden Punkte verwendet werden:
\begin{align*}
\text{P}_1 \left( \begin{array}{c} 2\\ 2\end{array} \right),
\text{P}_2 \left( \begin{array}{c} 3\\ 3\end{array} \right),
\text{P}_3 \left( \begin{array}{c} 1\\ 4\end{array} \right),
\text{P}_4 \left( \begin{array}{c} 5\\ 3\end{array} \right),
\text{P}_5 \left( \begin{array}{c} 5\\ 7\end{array} \right)
\end{align*}
\end{document}
\clearpage
\section{Lösungsvorschlag}

\begin{figure}[h!]
  \begin{tikzpicture}
      \draw[step=.5 cm, gray!90, very thin] (0,0) grid (6,3);
      \draw[->] (-.5,0) -- (7,0);
      \draw[->] (0,-.5) -- (0,4);
      %\draw node[circ] at (0,0){};
      \filldraw[white] (1,1) circle [radius=5pt]
                       (2,2) circle [radius=5pt]
                       (5,2) circle [radius=5pt];
      \draw (1,1) circle [radius=5pt]
            (2,2) circle [radius=5pt]
            (5,2) circle [radius=5pt];
      \draw node at (1,1) {+};
      \draw node at (2,2) {+};
      \draw node at (5,1.95) {-};
      %\draw node at (3,.95) {-};
      \draw node at (0.5,1) {$Q_1$};
      \draw node at (5.5,2) {$Q_2$};
      \draw node at (1.5,2) {$P_1$};
%       \draw[<->] (1.2,1.2) -- (1.8,1.8);
    %  \draw[>-<] (2.8,1.2) -- (2.2,1.8);
%       \draw[->] (2,2.3) -- (2,3);
%       \draw[->] (2.3,2) -- (3,2);
  \end{tikzpicture}
  \caption{Abstände Q1, Q2 zu P1 (Ein Kästchen = 0,5 m)}
\end{figure}

\subsection{Abstand}
\subsubsection{Abstand Q1 - P1}
\begin{align}
  |P_1-Q_1| &= \sqrt{(P_{1x}-Q_{1x})^2 + (P_{1y}-Q_{1y})^2}\\
  |P_1-Q_1| &= \sqrt{(2-1)^2 + (2-1)^2}\\
  &=\sqrt{2}\, m \approx 1,41\, m
\end{align}

\subsubsection{Abstand Q2 - P1}
\begin{align}
  |P_1-Q_2| &= \sqrt{(P_{1x}-Q_{2x})^2 + (P_{1y}-Q_{2y})^2}\\
  |P_1-Q_2| &= \sqrt{(2-5)^2 + (2-2)^2}\\
  |P_1-Q_2| &= \sqrt{3^2 + 0^2}\\
  &=3\, m
\end{align}

\subsection{Elektrisches Feld}
\begin{align}
\vec{E}_{Pn} &= \frac{Q_n}{4\pi \varepsilon_0 r^2}\\
\vec{E}_{Q1P1} &= \frac{1,2 nC}{4\pi \varepsilon_0 \sqrt{2}}\\
   &= 7,63\ \frac{V}{m}\\
\vec{E}_{Q2P1} &= \frac{-15\ nC}{4\pi \varepsilon_0 3}\\
 &= -44,94\ \frac{V}{m}\\
\end{align}
\subsubsection{resultierendes Feld}
\begin{align}
 \vec{E}_{ges} &= \vec{E}_{Q1P1} + \vec{E}_{Q2P1}\\
  &= 7,63\frac{V}{m} -44,94\ \frac{V}{m}\\
  &= -37,31\ \frac{V}{m}
\end{align}

\subsection{Kraft auf Q3 am Punkt P1}
\begin{align}
  \vec{F}_{QnPn} &= \frac{Q_n\cdot Q_{P_n}}{4\pi \varepsilon_0 r^2}\\
\vec{F}_{Q1P1} &= \frac{1,2\ nC\cdot 33\ nC}{4\pi \varepsilon_0 \sqrt{2}}\\
   &= 251,66\, nN\\
\vec{E}_{Q2P1} &= \frac{-15\ nC\cdot 33\ nC}{4\pi \varepsilon_0 3}\\
 &= -1,48\ \mu N\\
\end{align}
\subsubsection{resultierende Kraft}
\begin{align}
 \vec{E}_{ges} &= \vec{E}_{Q1P1} + \vec{E}_{Q2P1}\\
  &= 251,66 \ nN -1,48\ \mu N\\
  &= -1,23\ \mu N
\end{align}

\subsection{Winkel Q1-P1}
$sin(45^\circ) = 0,71$\\
$cos(45^\circ) = 0,71$

\subsection{Winkel Q2-P1}
$sin(0^\circ) = 0$\\
$cos(0^\circ) = 1$


